% ============================================================================
\section{Justificación de la investigación}

\subsection{Justificación teórica}

La investigación contribuyó a la literatura de revisiones estadísticas al formular un problema predictivo explícito: anticipar el cambio entre una primera publicación y una publicación posterior definida mediante un horizonte temporal fijo. También propuso integrar la caracterización econométrica de las revisiones con una comparación entre modelos clásicos y una arquitectura QRC. La contribución no consistió en asumir una superioridad cuántica, sino en establecer condiciones bajo las cuales una arquitectura de reservorio podría o no aportar información incremental.

\subsection{Justificación metodológica}

El estudio estableció un protocolo de validación temporal que impidió utilizar información futura en el entrenamiento, la normalización, la imputación y la selección de hiperparámetros. Además, propuso calibrar el experimento sintético a parámetros estimados en las revisiones observadas, propagando la incertidumbre de calibración. La comparación incluyó controles clásicos y métricas predictivas complementarias, con inferencia adecuada a la dependencia temporal.

\subsection{Justificación práctica}

Una mejor caracterización de las revisiones ayuda a que los analistas interpreten con mayor prudencia los primeros anuncios del PBI. Si se identifica un patrón predecible, esa información puede servir como referencia complementaria en el seguimiento coyuntural. No obstante, el estudio no presupone que una mejora estadística produzca automáticamente mejores decisiones de política monetaria, pues cualquier recomendación institucional depende de la estabilidad, oportunidad y costo operativo del método empleado.

\subsection{Relevancia para el Perú}

La investigación es pertinente para el Perú porque interpretar correctamente la coyuntura exige considerar la calidad y las actualizaciones posteriores de las cifras oficiales. Una evaluación ordenada de publicaciones sucesivas documenta la incertidumbre inherente al dato preliminar y ofrece un marco reproducible para futuros estudios, incluso cuando los modelos cuánticos no superen a las alternativas clásicas convencionales.
